\documentclass[11pt,a4paper]{article}
\usepackage[T1]{fontenc}
\usepackage{lmodern}
\usepackage[margin=26mm]{geometry}
\usepackage{amsmath,amssymb,amsthm,booktabs,array}
\usepackage{microtype}
\usepackage[hidelinks]{hyperref}
\hypersetup{pdftitle={Deterministic Integer 3SUM at Exponent 1.960614},pdfauthor={Yuxuan Xu}}
\newtheorem{theorem}{Theorem}[section]
\newtheorem{lemma}[theorem]{Lemma}
\newtheorem{proposition}[theorem]{Proposition}
\newcommand{\R}{\mathbb{R}}
\newcommand{\N}{\mathbb{N}}
\newcommand{\Oh}{\mathcal{O}}
\newcommand{\code}[1]{\texttt{#1}}
\setlength{\emergencystretch}{3em}
\title{Deterministic Integer 3SUM at Exponent 1.960614\\[4pt]
\large Quartic Switch Certificates and a Lean Word-RAM Proof}
\author{Yuxuan Xu\\[3pt]
\small School of Computer Science, Wuhan University\\
\small ORCID: \href{https://orcid.org/0009-0001-8760-5277}{0009-0001-8760-5277}}
\date{October 15, 2026\quad---\quad Technical report, version 1.0.0}
\begin{document}
\maketitle
\begin{abstract}
This report documents a deterministic integer 3SUM upper bound of
$O(n^{1.960614})$ in the precise word-RAM model used by the accompanying
Lean development. It extends the pinned \code{c0ldzy17/three-sum-lean}
formalization with local expansion $c=232$, a rational quartic switch,
integer schedules, and revised execution and resource proofs. A uniform
entropy inequality follows from two supporting lines; exact rational
interval arithmetic certifies every entry of a twelve-term cost ledger
below $1.96061388$. The complete formal endpoint is
\code{EndStatement.ThreeSum.SolvedInTime (980307/500000 : Rat)}.
Archived official verification records for the pinned source report
successful Lean and Nanoda kernel checks, statement comparison, and an
allowed-axiom audit. This report explains the analytic certificates,
the connection to the formal artifact, and the boundaries of the evidence.
The full machine proof is supplied as source rather than reproduced here
as a standalone paper proof. The deposit is a technical report and proof
artifact, not a claim of journal peer review or a general 3SUM lower bound.
\end{abstract}
\noindent\textbf{Keywords:} 3SUM; deterministic algorithms; word RAM;
Lean; formal verification; entropy bounds; exact certificates.

\section{Statement and contribution}
The input consists of $n$ integers. The decision problem asks whether
three pairwise distinct positions have values summing to zero; repeated
values at different positions are allowed. The result concerns
polynomially bounded integer inputs, not arbitrary real inputs.

\begin{theorem}[Formal endpoint]\label{thm:end}
For each fixed $\kappa\in\N$, there exist a deterministic word-RAM
program $P$, a word-width coefficient $b\in\N$, and a time function
$T:\N\to\N$ with $T(n)=O(n^{980307/500000})$ such that $P$ correctly
decides 3SUM for every input with $|a_i|\le n^\kappa$, at every word width
$W\ge b(\lfloor\log_2 n\rfloor+1)$, within $T(n)$ steps.
\end{theorem}

The instruction set in the unchanged specification consists of the
constant one, addition, subtraction, multiplication, indirect load and
store, a branch on a negative signed word, and accept/reject instructions.
Each executed instruction costs one step. The statement quantifies one
program per fixed magnitude exponent, uniformly over all sufficiently
large word widths. The exact rational exponent is $980307/500000=1.960614$.

Theorem~\ref{thm:end} is the type of \code{Race.bound} in
\code{proof-source/Submission.lean}. Its proof uses
\code{ThreeSum.solvedInTime\_980307\_500000} in
\code{Submission/ThreeSum/FinalSolver.lean}. The exact operational
semantics and quantifiers are in \code{Submission/EndStatement.lean}.
For a rational exponent $r$, that specification expresses the time bound as
$\exists K\in\N\ \forall n\ge2,\ T(n)^{r.\mathrm{den}}
\le K n^{r.\mathrm{num}}$.

The construction inherits the algorithmic framework, reductions, machine
model, and substantial proofs from earlier work \cite{avw,base,anthropic}.
The changes documented here are the expansion-232 instantiation, a
positive-coefficient quartic certificate, rounded integer schedules,
and the machine-level proof changes needed to connect them to the exact
endpoint. The baseline is pinned to commit
\code{2b639655e07421b9ac743de9090459bab473a7c0}; the verified extension is
pinned to \code{8c25c0325de5be4286413f39ca0c9278aad33f52} \cite{artifact}.
No exhaustive novelty search or claim to invent the inherited framework
is made. The earlier manuscript at exponent $1.999017$ concerns a
different stage of the analysis and is not relabelled as this result.

\section{Entropy functions and an exact quartic switch}
Let $H(t)=-t\log_2t-(1-t)\log_2(1-t)$, continuously extended at
$0$ and $1$. For an integer $c\ge3$, set
\begin{align}
 E_c&=\tfrac c2 H(1/c),\\
 k_c&=2cH(1/(2c))-cH(1/c)-1,\\
 \gamma_c(t)&=cH(1/c)-cH((1-t)/c).
\end{align}
The notation $k_c$ denotes a local combinatorial exponent and is distinct
from the input-magnitude exponent $\kappa$.

\begin{lemma}[Positive expansion]\label{lem:series}
For $0\le t\le1/2$,
\[
 \gamma_c(t)=\log_2(c-1)t+
 \frac1{\ln2}\sum_{j=2}^{\infty}
 \frac{1+(-1)^j/(c-1)^{j-1}}{j(j-1)}t^j.
\]
All coefficients in this expansion are positive.
\end{lemma}
\begin{proof}
Differentiate to obtain
\[
 \gamma_c'(t)=\frac{\ln(c-1)+\ln(1+t/(c-1))-\ln(1-t)}{\ln2}.
\]
The logarithm series converge absolutely and uniformly on the interval.
Integrating term by term from zero gives the formula. Since $c-1\ge2$,
the displayed coefficient numerator is strictly positive for every $j$.
\end{proof}

For $c=232$, use the polynomial
\begin{equation}\label{eq:quartic}
 P(t)=7.85174904t+0.72447023t^2+0.24044466t^3+0.12022459t^4.
\end{equation}
Every decimal here and below is an exact rational. Rational logarithm
enclosures verify that each coefficient is below the corresponding
coefficient of Lemma~\ref{lem:series}. Thus $P\le\gamma_{232}$ on
$[0,1/2]$. Positivity of the coefficients gives monotonicity and convexity;
the same certificate verifies $P'(1/2)<9$.

\section{A uniform entropy inequality}
The switch requires an inequality valid on an entire interval, not just
at sampled values. Fix
\[
 \rho=0.27905030,\qquad u=0.17200068185091493,
\]
and define $A=-\log_2(1-u)$ and $B=\log_2((1-u)/u)$.
Concavity of entropy and convexity of $P$ give
\[
 H(t)\le A+Bt,\qquad
 P(t)\ge P(u)+(t-u)P'(u).
\]
The difference
\[
 L(t)=\rho[1+P(u)+(t-u)P'(u)]-(A+Bt)
\]
is affine. The exact certificate supplies the strict bounds
\[
 L(0)>1.486795215\cdot10^{-9},\qquad
 L(1/2)>7.676481609\cdot10^{-9}.
\]
Linear interpolation proves $L(t)>0$ throughout $[0,1/2]$, including its
endpoints. Consequently,
\begin{equation}\label{eq:uniform}
 H(t)<\rho(1+P(t))\qquad(0\le t\le1/2).
\end{equation}
The interval bounds certify the two endpoint comparisons; the supporting
line argument explains why those comparisons imply a continuous theorem.

\section{Finite schedules and local expansion}
Use the parameters
\begin{equation}\label{eq:params}
\begin{split}
 e&=0.203889696,\qquad b=-0.062656016,\\
 S&=0.193175570,\qquad \theta=0.040409899,\qquad \tau=13/20,
\end{split}
\end{equation}
and put $v=S-2b>0$. Here $b$ is the band-shift parameter, not the
word-width coefficient of Theorem~\ref{thm:end}.

\subsection{Integer sizes and the switch comparison}
Let $m=\lfloor\log_2 n\rfloor$ and take the even microblock size
$s=2\lfloor\lfloor m^{1/3}\rfloor/2\rfloor$. For $s>0$, define
\[
 J=20\left\lfloor\frac{em}{20s}\right\rfloor,\qquad \mu=Js.
\]
Then $20\mid J$ and $0\le em-\mu<20s$, so $\mu/m\to e$.
Use the integer band divisor $2^{\lfloor-b\mu\rfloor}$. It differs from
the ideal real value by a factor below two. The divisibility of $J$
ensures the terminal depth $\tau J$ is integral. Exceptional small sizes
are handled by the finite fallback branch of the formal program.

Write $P(t)=\sum_{i=1}^4 a_i t^i/Q$ with $Q=10^8$, and $v=N/D$.
At depth $j\le\tau J$, let $a=j/J$ and $r=(J-j)s>0$. Select the least
integer $q\in[0,r/2]$ satisfying
\begin{equation}\label{eq:integer}
 D(J-j)\sum_{i=1}^4 a_iq^ir^{4-i}\ \ge\ Q(NJ+Dj)r^4.
\end{equation}
All denominators are positive, so this comparison is exactly equivalent to
$P(q/r)\ge(v+a)/(1-a)$. It uses only integer arithmetic. The certificate
checks $P(1/2)>(v+\tau)/(1-\tau)$; since $r$ is even, an admissible $q$
exists at every such depth. Minimality of $q$, together with $P'<9$, gives
\[
 \frac{v+a}{1-a}\le P(q/r)<\frac{v+a}{1-a}+\frac9r.
\]
With $t=q/r$, this implies
\begin{align}
 S&\le2b-a+(1-a)\gamma_{232}(t),\\
 (1-a)H(t)&<\rho(1+v)+9\rho/\mu.
\end{align}
The certificate also checks $P(\theta)>v$. The initial integer threshold
is at most $\theta+1/\mu$ for the large-size branch. Its entropy is bounded
using the tangent to $H$ at $\theta$. These facts provide uniform rounding
control. The machine proof additionally accounts for scalar bounds,
memory layout, and the execution cost of these comparisons.

\subsection{Why the expansion can be 232}
For $s=2h>0$ and $L=cs$, let
\[
 M=\binom Ls,\quad V=\binom Lh,\quad
 \Delta=\binom{L-h}h,\quad G=\binom sh.
\]
Double counting disjoint ordered $h$-subsets gives $V\Delta=MG$.
Since $\Delta\le V$ and $G<2^s$, we obtain $\Delta^2<M2^s$.
Choosing one element from each of $s$ disjoint $c$-element sets also gives
$M\ge c^s>2^s$. A prime selected above
$\lceil\sqrt{M2^s}\rceil$ and within twice that integer satisfies
$M2^s\le p^2\le16M2^s$, hence $\Delta<p$ and $G<p$.
These finite inequalities do not require $c$ to be a power of two.

Standard binomial entropy estimates yield
\[
 \log_2 M=cH(1/c)s+O(\log s),\qquad
 \log_2 V=cH(1/(2c))s+O(\log s),
\]
and $\log_2p=(E_c+1/2)s+O(\log s)$. Moreover, $s^2/m\to0$ and
$J\log(s)/m\to0$, so the associated local enumeration and polynomial
factors are $n^{o(1)}$. These estimates alone do not certify the whole
program: the concrete dimensions, bucket construction, reporting routines,
and resource contracts must all use the same expansion. The accompanying
extension supplies that machine-level connection.

\section{The exact twelve-term certificate}
The cost ledger retained from the baseline has the following exponent
expressions, evaluated at $c=232$ and the parameters above:
\begin{center}\small
\begin{tabular}{@{}rl@{}}
\toprule
Term & Exponent expression\\\midrule
1 & $2eE_c$\\
2 & $1+e(E_c-b)$\\
3 & $2-eS$\\
4 & $2-e(1/2+b-H(\theta))$\\
5 & $2-e[1/2-b-k_c-\rho(1+S-2b)]$\\
6 & $2-e(\tau-1/2-b-k_c)$\\
7 & $\frac85e(E_c+1/2)$\\
8 & $2+e(4/5-\frac25E_c-2b)$\\
9 & $1+e(1/2-b)$\\
10 & $2-e(E_c+b)$\\
11 & $2-2e(E_c+b)$\\
12 & $1$\\\bottomrule
\end{tabular}
\end{center}
Exact interval evaluation proves
\begin{align*}
4.648781151243719&<E_{232}<4.648781151243720,\\
0.001557987216578&<k_{232}<0.001557987216579,\\
H(\theta)&<0.244168414687344.
\end{align*}
The largest certified upper endpoint is
\[
 1.9606138795768596941522816<1.96061388<1.960614.
\]
The gap is positive and therefore absorbs the proved $n^{o(1)}$ overhead
for sufficiently large inputs; finite exceptions are handled by the
formal fallback. When $E_c$ appears with a negative coefficient, its lower
endpoint must be used to obtain an upper cost bound. The checker enforces
this sign-sensitive interval calculation.

For completeness, the archived analytic package also includes degree-eight
certificates at $c=232$ and $c=256$. The formal endpoint presented here
uses the degree-four $c=232$ choice. A separate fixed-$c$, uniform-entropy
cost-model argument gives
\[
1.96061349120<\inf R_{\mathrm{model}}(232)<1.960613492.
\]
Its admissible parameters and proof are defined in Section 7 of the archived
notes:
\begin{center}\code{analytic-certificate/PROOF.md}.\end{center}
It is independent of the
algorithmic endpoint and is not a lower bound for arbitrary 3SUM algorithms.
It does not exclude changing the model, mixing expansions, or using
different algorithmic constructions.

\section{Connection to the complete formal program}
The verified source includes the concrete expansion-232 cover and bucket
construction, exact schedules, stored reporting data, host execution and
resource bounds, and the final hybrid solver. In particular, the last
theorem is not obtained by treating the numerical ledger as a machine
runtime theorem. \code{FinalSolver.lean} invokes
\code{ThreeSumHybridEndBridge.endStatement\_of\_solvers}, supplying the
actual hybrid program, its correctness proof, and its polynomial resource
bound for every magnitude exponent. The compiler bridge then yields the
original word-RAM statement. The source is the authoritative detailed
proof of this composition.

The archived official verifier receipt reports that \code{Race.bound}
passed both the Lean and Nanoda kernels and that its statement matched
the official frozen problem. Its reported axioms are exactly
\code{propext}, \code{Classical.choice}, and \code{Quot.sound}, with no
\code{sorryAx}. The receipt is dated October 14, 2026, and is tied to the
source commit above and verifier commit
\code{51489ee2ac82596a41375e92ea8108814f931242} \cite{verifier}.

The separately archived Prove2Me completion receipt, dated October 15,
records completion of all seven original milestones, 41 published machine
definition bundles, 481 accepted machine proof records, and 110 reused
proved theorem nodes. The 481 records correspond to 480 new machine theorem
declarations, not 481 independently novel mathematical results. A graph
audit records 1,316 actual dependency paths. These are platform provenance
statistics; they are not measures of novelty or substitutes for the
standalone native source and verifier checks.

The official audit also records \code{site\_eligible\_for\_record: false}.
This report records successful verification without interpreting a
leaderboard position as eligibility for a new speedrun record or as an
independent assessment of scientific priority.

\section{Reproduction and contents of the deposit}
The source archive accompanying this report contains the pinned native
proof source, the exact analytic certificates and their checker, selected
verification evidence, original license notices, and a SHA-256 manifest.
It excludes working-directory caches, credentials, unrelated research,
and the incomplete later exponent-improvement work.

The analytic checker needs Python 3.10 or newer and only the standard
library. From the extracted archive, run
\begin{verbatim}
python verify_release.py
python analytic-certificate/verify_certificate.py --self-test
\end{verbatim}
The first command checks archived bytes and the official Lean-source hash
map. It is an integrity check, not a new theorem proof. The second checks
three exact certificates and twelve deliberately corrupted negative cases.
The negative cases test rejection of wrong entropy ratios, oversized
polynomial coefficients, overly small cost caps, and overly large
model lower bounds.

Logarithms are enclosed using rational arithmetic. Reduce a positive
rational argument to $y\in[1,2]$ by powers of two and put
$z=(y-1)/(y+1)$. For $K=48$,
\[
 S_K\le\ln y\le S_K+
 \frac{2z^{2K+1}}{(2K+1)(1-z^2)},\qquad
 S_K=2\sum_{i=0}^{K-1}\frac{z^{2i+1}}{2i+1}.
\]
The omitted series has positive terms, and the displayed geometric bound
controls its remainder. Interval arithmetic accounts for negative range
reduction exponents. The continuous lemmas in this report explain the
logical use of the resulting finite inequalities.

Full formal reproduction uses Lean \code{v4.33.1}, with Mathlib pinned to
\begin{center}\code{0df444a360eaa60ab8c11dca51a86af692955474}.\end{center}
It uses the pinned official verifier environment.
The source follows the verifier's submission layout;
the verifier supplies the Lake configuration. A root-directory build of
the broader research workspace is not a reproduction of this theorem.
The deposit preserves the dependency manifest and official logs. Preparing
this deposit does not constitute a fresh full Lean or Nanoda run; the
kernel conclusions above are attributed to the archived official run.

The original analytic notes retain their October 14 statement that the
machine bridge was not yet established by that certificate package. Those
historical files are preserved byte for byte. The subsequent native
machine proof and official records are separate evidence and must not be
confused with what the arithmetic checker itself proves.

\section{Attribution and use of AI}
The author is responsible for the submitted record. OpenAI Codex assisted
the extension and formalization, inspection and organization of the
artifacts, and preparation of this report. Automated checks and platform
acceptance are not described as independent human peer review. This report
does not assert that a separate human line-by-line audit has been completed.

Attribution to Alman and Vassilevska Williams, the baseline formalization,
Anthropic, and other upstream contributors is retained. In particular,
the inherited \code{LICENSE}, \code{NOTICE}, and per-file copyright
headers accompany the source unchanged. The baseline's acknowledgments
include Bj\"orklund, Husfeldt, Kaski, Koivisto, Swapnil Jain, and the
formalization lineage through Jihoon Hyun. The deposit's licensing note
distinguishes the report from inherited software; it does not reassign
ownership or license third-party contributions on the author's behalf.

\begin{thebibliography}{9}\raggedright
\bibitem{avw} Josh Alman and Virginia Vassilevska Williams.
\emph{Truly Subquadratic 3SUM and Truly Subcubic APSP via Triangles in Sparse
Lopsided Graphs}. arXiv:2610.06783v1, 2026.
\url{https://arxiv.org/abs/2610.06783v1}.
\bibitem{base} c0ldzy17. \emph{three-sum-lean}, pinned source and METHOD.md,
commit \code{2b639655e07421b9ac743de9090459bab473a7c0}, 2026.
\url{https://github.com/c0ldzy17/three-sum-lean/tree/2b639655e07421b9ac743de9090459bab473a7c0}.
\bibitem{anthropic} Anthropic. \emph{Truly Subquadratic 3SUM and Truly
Subcubic APSP: a Lean 4 formalization}. Source lineage and copyright
notices are preserved in the accompanying \code{proof-source/NOTICE}
and \code{proof-source/Submission/} tree.
\bibitem{artifact} Yuxuan Xu (Anchen0823).
\emph{three-sum-1960614}, commit
\code{8c25c0325de5be4286413f39ca0c9278aad33f52}, 2026.
\url{https://github.com/Anchen0823/three-sum-1960614/tree/8c25c0325de5be4286413f39ca0c9278aad33f52}.
\bibitem{verifier} Big O Speedrun. Official verification issue and run
for the preceding source, October 14, 2026.
\url{https://github.com/bigospeedrun/verifier/issues/8};
\url{https://github.com/bigospeedrun/verifier/actions/runs/38567645711}.
\bibitem{platform} Prove2Me. \emph{Deterministic Integer 3SUM 1.960614:
Quartic Switch and Complete Program Extension}. Completion receipt
archived October 15, 2026; campaign
\url{https://prove2.me/campaigns/3sum-exponent}.
\end{thebibliography}
\end{document}
